\section{Related Literature}

Eye tracking is being adopted heavily in XR, providing clear value in human-computer interaction (HCI) interfaces.  As contextual AI emerges, new prototypes have explored eye gaze as a means to estimate user attention.  This section provides an overview of related literature, including the use of ET for selection, scene understanding, and ET in contextual AI.

\subsection{Eye Tracking for Selection in Extended Reality}

In recent years, eye gaze has gained popularity as a signal for HCI in XR systems~\cite{plopski_survey_2022}.  Eye gaze has comparable usability to controllers~\cite{luro_comparative_2019, zhang_telepresence_2019,fernandes2024effect} while freeing the hands for other tasks and being preferable to users~\cite{piening_looking_2021}.  While ET is prone to a \textit{Midas touch} fallacy, where false selections are made during ambient fixations~\cite{jacob_advanced_1995}, novel HCI methodologies~\cite{khamis_vrpursuits_2018} overcome this and make ET an ideal signal for interface navigation. Our work explores ET's ability to continuously and implicitly convey a user's attention and intent~\cite{sendhilnathan2024implicitgazeresearchxr}.  This has only lightly been explored with AI agents, though ET has facilitated automatic contextual displays in the past~\cite{toyama_recognition_2012}.

\subsection{Eye Gaze Encodes Scene Understanding}

Eye movements reflect a viewer's internal processing of a scene, giving insights to cognitive state and attention as one interprets new visual stimuli~\cite{langton_eyes_2000, eckstein_beyond_2017}. Sequences of gaze fixations (i.e., scanpaths) encodes contextual cues as to future objects of interest~\cite{itti_computational_2001, burlingham_laziness_2024}; a number of works have leveraged scanpath history for short-term gaze prediction / anticipation~\cite{huang_predicting_2018, david-john_prediction_2021, Hu2021, damelio_tppgaze_2024}. Burlingham et al. found temporal dependencies in scanpaths lasting for 4-5 fixations on average, with high variance across task types~\cite{burlingham_scanpaths_2024}. Contextual AI models may be able to leverage the rich, multiscale structure of scanpaths when inferring intent. %(analogous to how large language models (LLMs) leverage multiscale structures in language), enabling implicit inferences about user intent from scanpath history.

Insights about the cognitive encoding of nearby objects can inform our expectations for eye movements in contextual AI~\cite{Tatler2011}.  For example, humans tend to look at a coffee mug just before grasping, to encode the location of the handle so that it can be successfully grasped.  Object locations are encoded via an egocentric reference to the user, and greater affordances are given to nearby interactable objects~\cite{costantini_affordance_2010, Tatler2011}.  The visual system elicits responses to reachable 3D objects, even when there is no intent to interact~\cite{iachini_motor_2014, iachini_temporal_2023}.  So, by analyzing the eye gaze fixations on nearby objects, we could predict possible future interactions. 

\subsection{Eye Tracking in Contextual AI}

Information from the physical world could greatly improve user interactions with AI agents~\cite{zhang_vision_survey_2024}.  Emerging products, such as the Ray-Ban Meta\footnote{\url{https://www.meta.com/smart-glasses}} and Google's Ask Photos\footnote{\url{https://blog.google/products/photos/ask-photos-google-io-2024/}}, use image context to improve user interaction.  But, given a full image, the human's intent may not align with the salient features detected by the agent.  Eye gaze could aid in narrowing relevant context observed by contextual agents~\cite{ajanki_contextual_2010, buschel_here_2018}, enhancing understanding and avoiding hallucination~\cite{cui_holistic_2023, leng_hallucinations_2024}.  

Some wearable contextual AI prototypes integrating eye tracking have been proposed~\cite{zhang_agi_2024}.  The GazeGPT system projects 2D gaze onto an image capture, cropping image contents before interfacing with a VLM~\cite{konrad_gazegpt_2024}.  They demonstrated gaze-based querying to be faster, more accurate, and more natural than head-mounted and smart-phone-like baselines.  G-VOILA interfaces with a textual LLM, using gaze-generated saliency maps for prompt adjustment~\cite{wang_gvoila_2024}.  Derived object information is spliced into the query, increasing robustness against ambiguity and increasing participants' confidence in the system.  These prototypes show clear value from the inclusion of ET for point-in-time querying.  